\section{Текстовые спецификации основных классов}
\subsection{Общие соглашения интерфейсов}
Имена и сигнатуры ниже соответствуют заголовкам в \file{src/core} и \file{src/ui}. Возвращаемый \code{CalendarResult} содержит код \code{CalendarError}, пояснение, ID мероприятия, конфликтующие ID и ID заявки; \code{ok()} означает отсутствие ошибки. Отказ не подтверждает конфликтующее мероприятие. Некорректные настройки конструктора вызывают исключение \code{std::invalid_argument}; ошибки файлов и воспроизведения --- \code{std::runtime_error}.

\code{EventId} --- беззнаковый 64-битный идентификатор. Данные передаются преимущественно по константной ссылке; необязательные ссылки представлены \code{std::optional}. Все операции с модельным временем используют \code{QDateTime}.

\subsection{EventData и CalendarEvent}
\textbf{Назначение.} Представление пользовательских полей и управляемого сервисом состояния мероприятия. \code{EventData} хранит название, тип, начало и конец, необязательное помещение, ID участников, важность, организатора и необязательный ID шаблона повторения. \code{CalendarEvent} содержит \code{EventData}, ID, статус, две версии и снимки кодов и отделов участников.

\textbf{Инварианты при принятии.} Непустое название; существующие участники без повторов; непустой состав; существующий организатор; корректный интервал с концом после начала; допустимые перечисления; существующее помещение при его указании; отсутствие помещения у командировки. Начало не раньше текущего времени. \code{version} меняется при редактировании, \code{scheduleVersion} --- только при изменении времени, важности или состава участников. Клиент не задаёт ID и статус создаваемого события напрямую.

\subsection{CalendarSetup}
\textbf{Назначение.} Совокупность параметров одного опыта: каталог компании, период и шаг, рабочее время, генератор, начальные и периодические планы, политика конфликтов, правила напоминаний, роль и количество повторов. Структура не продвигает время и не содержит виджетов.

\textbf{Проверка.} \code{makeCompanyCatalog(setup)} создаёт каталог; конструктор движка проверяет согласованность параметров. Допустимы 5--9 отделов, 7--30 дней, шаг 30/60 минут, 1--1000 сотрудников на отдел. Рабочая неделя должна содержать хотя бы один рабочий день; начало рабочего времени раньше конца. Число участников не превышает штат. Интенсивности неотрицательны; веса важности имеют положительную сумму. Для включённого начального пресета требуется рабочий день минимум один час.

\subsection{CalendarService}
\textbf{Назначение и состояние.} Единственный владелец активного календаря, архива, истории и каталога одного прогона. Хранит текущий момент и конец горизонта. Не генерирует случайные события и не рисует интерфейс.
\contract{CalendarService(catalog, start, horizonEnd)}{Каталог и корректные границы времени.}{Инициализированный пустой календарь.}{Начало строго раньше конца; ссылки каталога согласованы; коды сотрудников уникальны. Некорректная конфигурация вызывает исключение.}
\contract{checkEvent(data, actor, id=0)}{Предлагаемые данные, контекст пользователя, необязательный ID редактируемого плана.}{Результат проверки и список конфликтующих событий.}{Не изменяет состояние. Для id=0 проверяет создание, иначе существование, полномочия и статус редактируемого события.}
\contract{createEvent(data, actor)}{Поля нового события и контекст пользователя.}{Результат с новым EventId при успехе.}{Выполняет общую проверку; присваивает ID и версии, фиксирует снимки и запись Created. При допустимом начале в текущий момент событие может стартовать сразу, если старты не отложены движком.}
\contract{updateEvent(id, data, actor)}{ID и новые поля существующего плана.}{Успех либо причина отказа.}{Изменяется только Planned; проверяются права на исходные и новые данные. При успехе обновляются версии и история. Отказ сохраняет прежнее событие.}
\contract{calendar(filter)}{Фильтр области, дат, часов, отдела, сотрудника, помещения и статуса.}{CalendarQueryResult со списком и результатом проверки фильтра.}{Отбор по пересечению интервалов; просмотра достаточно для фильтрации, состояние не меняется.}
\contract{advanceTo(time, startAtBoundary=true)}{Неубывающий момент внутри горизонта и флаг обработки стартов.}{Результат продвижения времени.}{Обрабатывает все начала и окончания внутри интервала; завершения переносятся в архив. Режим false позволяет движку обработать заявки и уведомления до граничных стартов.}
\contract{cancelEvent(id, actor)}{ID запланированного события и контекст.}{Успех либо отказ.}{Проверяет полномочия и статус, переносит событие в архив как Cancelled. Метод есть в ядре, но отдельной кнопки отмены в UI нет.}
\textbf{Читающие операции.} \code{event(id)} возвращает необязательное событие; \code{history(id)} --- историю всех или одного события; \code{historySince(offset)} --- её хвост; \code{participantCode} и \code{participantKey} преобразуют внутренние и публичные ID. \code{employeeBusy} проверяет участие в текущем мероприятии.

\subsection{SimulationEngine}
\textbf{Назначение.} Координация одного опыта. Содержит календарь, параметры, генератор \code{std::mt19937}, напоминания, почтовый канал, заявки, журнал и команды. От Qt Widgets не зависит.
\contract{SimulationEngine(setup)}{Параметры CalendarSetup.}{Готовый к первому шагу движок.}{Проверяет параметры, создаёт каталог и исходные регулярные планы при включённом пресете. Версия алгоритма --- calendar-v7.}
\contract{step()}{Явных параметров нет; используется сохранённый шаг.}{CalendarResult.}{Продвигает время на 30/60 минут с обработкой внутренних минут, не переходя конец периода. После конца возвращает отказ. Одна реализация используется таймером, кнопкой шага и ExperimentRunner.}
\contract{reset()}{Нет.}{Исходное состояние с теми же параметрами.}{Очищает историю, заявки, уведомления и команды, восстанавливает seed и начальный календарь. Это новый опыт, а не возврат к сохранённой точке.}
\contract{createEvent / updateEvent}{Данные мероприятия, контекст пользователя, при изменении ID.}{Результат с requestId.}{Регистрирует команду и заявку, обращается к сервису, применяет выбранную политику конфликтов, публикует историю и обновляет план напоминаний.}
\contract{acceptSuggestion(requestId, index)}{ID ожидающей заявки и индекс варианта с нуля.}{Результат новой попытки.}{Исходная заявка получает признак selected; выбранные данные повторно проходят общую проверку. Устаревший вариант не считается подтверждённым. Полномочия оператора в UI дополнительно проверяет RunTools.}
\contract{rejectSuggestion(requestId)}{ID ожидающей заявки.}{Результат снятия ожидания.}{Не отменяет календарное мероприятие; завершает рассмотрение неподтверждённой заявки. Используется при возврате к редактированию.}
\contract{trigger(kind)}{Тип календарного сценария.}{Результат сценарной заявки или переноса.}{Поддерживает пять типов мероприятий и изменение расписания. Некорректный тип и отсутствие подходящих планов приводят к отказу.}
\contract{replayCommand(command)}{Команда JSON с операцией, аргументами и временем.}{Изменение восстановленного движка; при ошибке исключение.}{Время команды должно совпасть с текущим; неизвестная операция отклоняется. Читающие методы возвращают календарь, заявки, письма, журнал и команды без права изменения.}

\subsection{ConflictResolver}
\textbf{Назначение.} Поиск вариантов без изменения календаря.
\contract{alternatives(calendar, setup, requested, actor, id=0)}{Текущий календарь, параметры поиска, исходная заявка, права и ID изменяемого события.}{QVector<EventData> размером не больше maxSuggestions.}{Сохраняет длительность и участников, перебирает рабочие даты, начала и помещения. Не принимает конфликтующий кандидат. Пустой результат означает отсутствие найденного варианта в заданном окне, а не невозможность любого расписания вообще.}

\subsection{ReminderService и канал уведомлений}
\textbf{Состояние.} \code{ReminderService} хранит для EventId версию расписания, первый и следующий сроки и период в секундах.
\contract{observe(records, rules)}{Новые записи истории и три правила по важности.}{Обновлённый план уведомлений.}{Создаёт или перестраивает серию при новой версии; прекращает её после старта, завершения или отмены. Та же версия повторно не планируется.}
\contract{deliver(calendar, channel)}{Текущий календарь и канал доставки.}{Вызовы send для наступивших сроков каждому участнику.}{Уведомляет строго до начала мероприятия и до конца периода; передвигает курсор, предотвращая дубли. Позднее создание компенсирует пропущенное окно одной отправкой.}
\contract{NotificationChannel::send(message)}{Значение EmailMessage с адресатом, событием, версией, текстом и временем.}{Контракт виртуальной операции доставки.}{Абстрактный метод; планировщик не зависит от конкретного транспорта. Виртуальный деструктор обеспечивает корректное уничтожение через базовый тип.}
\contract{SimulatedEmailChannel::send(message)}{Сформированное уведомление.}{Запись в messages().}{deliveredAt приравнивается к sentAt, сообщение добавляется в журнал. Сеть не используется; прежние тексты сохраняются как снимки.}

\subsection{SimulationReport}
\contract{create(engine)}{Движок в любом текущем состоянии.}{QJsonObject полного отчёта.}{Не изменяет модель; содержит конфигурацию, версию, события всех статусов, историю, заявки, письма, журнал, сотрудников и отделы. До завершения отчёт промежуточный.}
\contract{metrics(engine)}{Текущий движок.}{QJsonObject счётчиков и занятости.}{Учитывает фактически начатые мероприятия; обрезает занятость прошедшим временем. Будущие бронирования в занятость не входят.}
\contract{csv(report) / summary(metrics)}{Полный JSON-отчёт либо метрики.}{CSV UTF-8 либо текстовая сводка.}{В общем CSV столбцы section,key,value; вложенные объекты сериализуются и экранируются. Плоский CSV серии опытов создаёт ExperimentRunner.}

\subsection{SimulationStore и JsonCodec}
\contract{snapshot / save}{Движок; при сохранении путь файла.}{JSON состояния либо атомарно записанный файл.}{Содержит параметры, команды, контрольный календарь и SHA-256 отчёта вместе с состоянием генератора. Поддерживается формат 3 и только calendar-v7.}
\contract{restore / load}{JSON состояния либо путь к нему.}{\code{std::unique_ptr<SimulationEngine>}.}{Проверяет формат и версию, создаёт новый движок, проигрывает команды и сравнивает контрольные данные. Ошибка не подменяет открытый прогон. Состояния прежних генераторов отклоняются.}
\contract{saveConfiguration / loadConfiguration}{Параметры и путь либо путь к JSON настроек.}{Файл либо CalendarSetup.}{Параметры проверяются созданием движка; файл настроек не содержит накопленной истории. JsonCodec выполняет преобразования, проверки типов, дат и атомарную запись через QSaveFile.}

\subsection{ExperimentRunner}
\contract{ExperimentRunner(cases, repetitions)}{Непустой список сценариев и 1--1000 повторов.}{Готовая серия независимых опытов.}{Проверяет параметры сценариев и отсутствие переполнения seed+i.}
\contract{comparison(base)}{Базовые параметры.}{Шесть ExperimentCase.}{Три политики конфликтов при исходной и удвоенной интенсивности заявок. Остальные параметры сохраняются.}
\contract{advance()}{Нет.}{Признак оставшейся работы.}{Выполняет один шаг текущего движка; в пакетном Suggest после шага выбирает первый вариант. По завершении сохраняет метрики и переходит к следующему опыту.}
\contract{report() / csv()}{Текущее состояние серии.}{JSON с наблюдениями и агрегатами либо плоский CSV.}{Учитывает только завершённые опыты; незавершённая серия имеет complete=false. Для числовых метрик вычисляются среднее и выборочное стандартное отклонение.}

\subsection{Основные классы интерфейса}
\textbf{Widget : QWidget.} Проверяет поля параметров, открывает новое \code{SimulationWindow}, загружает настройки или состояние, открывает серию опытов. Экземпляр прогона получает собственную копию параметров.

\textbf{SimulationWindow : QWidget.} Владеет движком и таймером; показывает календарь, архив и журнал. \code{restoreState(engine)} заменяет движок проверенным состоянием. Приватные операции добавления, редактирования и выбора доступны через сигналы виджетов; при открытии формы время останавливается. \code{simulationEngine()} и \code{calendarService()} предоставляют константный доступ для проверок.

\textbf{RunTools : QWidget.} Отображает заявки, письма, общий список событий и отчёт. Перед выбором предложения проверяет полномочия оператора; затем движок повторно проверяет расписание. Через переданные обратные вызовы ставит прогон на паузу, обновляет представление и открывает форму заявки.

\textbf{OfficeAnimator : QObject.} \code{update(calendar,durationMs)} перестраивает расположение и перемещение графических объектов; \code{pause}, \code{resume}, \code{finish}, \code{retime} управляют визуальным движением. Сцена владеет графическими элементами, родительские QObject --- дочерними объектами Qt. Аниматор не меняет календарь.
